\begin{abstract}
Pixel intersection-over-union (IoU) is the wrong evaluation tool for linear features such as
roads: a small centerline offset or a width error collapses the score even when the extracted
road network is geometrically and topologically correct. We first introduce DTAF1, a
distance-tolerant, per-class F1 score built on tolerance-radius matching via Euclidean distance
transforms and class-agnostic across road, building, and point classes, and CBHM, a harmonic-mean
composite of centerline Dice (clDice) and boundary F1 deliberately designed as DTAF1's harsh foil.
Applying both to real SpaceNet and ISPRS Potsdam imagery surfaced two failure modes, each fixed
with an additive, non-breaking companion metric: \texttt{cbhm\_soft} (area-weighted,
non-collapsing) and DTAF1-Topo (a raster-skeleton APLS connectivity term closing DTAF1's
road-breakage blind spot). Building on this evaluation vocabulary, we then ask a training
question: can a single model learn to predict point (0D), linear (1D), and polygonal (2D) map
features \emph{jointly}, using a loss shaped by the same tolerance/topology/instance-matching
considerations that motivate these metrics, rather than a geometry-agnostic per-pixel loss? We
introduce \texttt{PLEMMultiTaskLoss}, a differentiable multi-task training loss combining three
geometry-aware terms --- a soft tolerance-band term (DTAF1 analog), a soft-clDice term (Shit
et~al.'s soft-skeletonization), and a CenterNet-style Gaussian-heatmap term standing in for
Hungarian-matched point F1, which has no tractable differentiable relaxation --- plus a masked
multi-source training scheme letting heterogeneous single-modality datasets (SpaceNet:
road+building; Potsdam: building+point) jointly supervise one 4-class model. We validate the loss
with 20 unit tests, synthetic per-term ablation sweeps confirming each term encodes the geometric
sensitivity it claims to, and a small real joint-training pipeline sanity check across 153 real
tiles, confirming all four loss terms learn jointly and that point-feature evaluation flows
end-to-end through the metric library on a real trained model's output.
\end{abstract}

\begin{CCSXML}
<ccs2012>
   <concept>
       <concept_id>10002951.10003317.10003338</concept_id>
       <concept_desc>Information systems~Geographic information systems</concept_desc>
       <concept_significance>500</concept_significance>
       </concept>
   <concept>
       <concept_id>10002951.10003317.10003347.10003350</concept_id>
       <concept_desc>Information systems~Evaluation of retrieval results</concept_desc>
       <concept_significance>500</concept_significance>
       </concept>
   <concept>
       <concept_id>10003752.10003809.10010031</concept_id>
       <concept_desc>Theory of computation~Computational geometry</concept_desc>
       <concept_significance>300</concept_significance>
       </concept>
   <concept>
       <concept_id>10010147.10010257.10010293.10010294</concept_id>
       <concept_desc>Computing methodologies~Neural networks</concept_desc>
       <concept_significance>500</concept_significance>
       </concept>
   <concept>
       <concept_id>10010147.10010257.10010293.10010294</concept_id>
       <concept_desc>Computing methodologies~Image segmentation</concept_desc>
       <concept_significance>500</concept_significance>
       </concept>
 </ccs2012>
\end{CCSXML}

\ccsdesc[500]{Information systems~Geographic information systems}
\ccsdesc[500]{Information systems~Evaluation of retrieval results}
\ccsdesc[300]{Theory of computation~Computational geometry}
\ccsdesc[500]{Computing methodologies~Neural networks}
\ccsdesc[500]{Computing methodologies~Image segmentation}

\keywords{geospatial segmentation evaluation, road extraction, building extraction,
tolerance-radius metrics, topology-aware metrics, centerline Dice, boundary F1, Average Path
Length Similarity, multiclass segmentation, differentiable loss functions, multi-task learning,
point/keypoint detection, soft-skeletonization}
